\section{Literal fixed-shift crosswalk}
\label{sec:literal-crosswalk}

Fix one nonzero integer \(h_0\).  At the high-\(\beta\) TPC-32
endpoint, retain the scales
\begin{equation}
 Q_X=X^{267/400+o(1)},\qquad
 J_X=X^{133/400+o(1)},\qquad
 C_X=\lfloor J_X\rfloor.
\label{eq:critical-scales}
\end{equation}
Let
\[
 A_X(n)=A_{C_X}(n)
\]
be the literal signed normalized-determinant coefficient after all
three matched channels and every native contribution in each
determinant bin have been aggregated \citep{WangTPC32}.  Define
\begin{equation}
 \Omega_X=\supp A_X,\qquad
 S_X=\#\Omega_X,\qquad
 Z_X=\sum_{n\in\Omega_X}A_X(n),\qquad
 D_X=\sum_{n\in\Omega_X}|A_X(n)|^2.
\label{eq:literal-statistics}
\end{equation}
TPC-32 gives the imported support envelope
\begin{equation}
 S_X\ll X^{o(1)}Q_X.
\label{eq:literal-support-envelope}
\end{equation}
When \(A_X\ne0\), view
\[
 a_X=(A_X(n))_{n\in\Omega_X}\in\C^{\Omega_X}
\]
and define
\begin{equation}
 \rho_X
 =
 \frac{|Z_X|}{\sqrt{S_XD_X}},
 \qquad
 \Flat_X
 =
 \frac{|Z_X|^2}{D_X}
 =
 S_X\rho_X^2.
\label{eq:literal-rho-flat}
\end{equation}

\begin{definition}[Literal same-frame surrogate]
\label{def:literal-surrogate}
A literal same-frame surrogate is a nonzero vector
\[
 b_X\in\C^{\Omega_X}
\]
whose coordinate \(b_X(n)\) corresponds to the same normalized
determinant bin \(n\) as \(A_X(n)\).  Define
\begin{equation}
 \delta_X
 =
 \frac{\|a_X-b_X\|_2}{\|b_X\|_2},
 \qquad
 \rho(b_X)
 =
 \frac{\left|\sum_{n\in\Omega_X}b_X(n)\right|}
 {\sqrt{S_X}\|b_X\|_2}.
\label{eq:literal-surrogate-statistics}
\end{equation}
No separately chosen surrogate support or channel-wise
renormalization is permitted in these definitions.
\end{definition}

\begin{theorem}[Literal surrogate stability certificate]
\label{thm:literal-surrogate-certificate}
Assume \(A_X\ne0\), \(S_X\ge2\), and \(\delta_X<1\).  Then
\begin{equation}
 |\rho_X-\rho(b_X)|\le\delta_X.
\label{eq:literal-rho-stability}
\end{equation}
In particular, if
\begin{equation}
 \rho(b_X)\le\frac{X^{o(1)}}{J_X},
 \qquad
 \delta_X\le\frac{X^{o(1)}}{J_X},
\label{eq:literal-surrogate-hypotheses}
\end{equation}
then
\begin{equation}
 \boxed{
 \rho_X\le\frac{X^{o(1)}}{J_X},
 \qquad
 \Flat_X\le X^{1/400+o(1)}.}
\label{eq:literal-surrogate-conclusion}
\end{equation}
\end{theorem}

\begin{proof}
Equation \eqref{eq:literal-rho-stability} is
\cref{cor:optimal-additive-stability} on the frame \(\Omega_X\).
The hypotheses \eqref{eq:literal-surrogate-hypotheses} therefore
give
\[
 \rho_X\le\frac{X^{o(1)}}{J_X}.
\]
Using \eqref{eq:literal-support-envelope},
\[
 \Flat_X
 =S_X\rho_X^2
 \le
 X^{o(1)}\frac{Q_X}{J_X^2}.
\]
The critical exponents in \eqref{eq:critical-scales} give
\[
 \frac{Q_X}{J_X^2}
 =
 X^{(267-266)/400+o(1)}
 =
 X^{1/400+o(1)}.
\]
\end{proof}

\begin{remark}[Logical status of the certificate]
\label{rem:certificate-status}
\Cref{thm:literal-surrogate-certificate} is a deterministic
implication.  It does not construct \(b_X\), bound \(\delta_X\), or
prove its zero-mode angle.  Verifying
\eqref{eq:literal-surrogate-hypotheses} for the growing literal
fixed-shift family is precisely additional arithmetic or structural
work.
\end{remark}

\begin{remark}[A numerical surrogate must resolve the small angle]
\label{rem:numerical-resolution}
A computation showing only
\(\delta_X=o(1)\) can accurately reproduce the determinant energy
while missing the critical zero-mode angle by many powers of
\(J_X\).  To serve as a rigorous certificate without further signed
error structure, it must control the global relative \(\ell^2\)
error at \(X^{o(1)}/J_X\) resolution on the literal frame.
\end{remark}

\subsection{Natural-support equivalence}

TPC-83 and TPC-88 identify a conditional natural mass--energy
branch in which
\begin{equation}
 S_X=X^{o(1)}Q_X.
\label{eq:natural-support}
\end{equation}
Here and below this denotes two-sided comparability up to subpower
factors \citep{WangTPC83,WangTPC88}.

\begin{proposition}[Critical angle equivalence on the natural branch]
\label{prop:natural-angle-equivalence}
Assume \eqref{eq:natural-support}.  Then
\begin{equation}
 \boxed{
 \Flat_X\le X^{1/400+o(1)}
 \quad\Longleftrightarrow\quad
 \rho_X\le\frac{X^{o(1)}}{J_X}.}
\label{eq:natural-angle-equivalence}
\end{equation}
Without the lower half of \eqref{eq:natural-support}, the angle
condition remains sufficient by
\eqref{eq:literal-support-envelope}, but the converse need not hold.
\end{proposition}

\begin{proof}
The identity \(\Flat_X=S_X\rho_X^2\) and
\eqref{eq:natural-support} give
\[
 \rho_X^2
 =
 X^{-267/400+o(1)}\Flat_X.
\]
Thus the flatness exponent \(1/400\) is equivalent to
\[
 \rho_X
 \le
 X^{-133/400+o(1)}
 =
 \frac{X^{o(1)}}{J_X}.
\]
The one-sided statement follows exactly as in
\cref{thm:literal-surrogate-certificate}.
\end{proof}

\begin{proposition}[Relative-\(\ell^1\) and principal-angle
thresholds]
\label{prop:canc-rho-thresholds}
On the natural post-bin mass branch of TPC-88,
\[
 \Seff(A_X)=X^{o(1)}Q_X,\qquad
 S_X=X^{o(1)}Q_X.
\]
Consequently,
\begin{equation}
 \rho_X=X^{o(1)}\canc_X.
\label{eq:canc-rho-natural}
\end{equation}
Hence the TPC-88 relative-cancellation threshold
\(\canc_X\le X^{o(1)}/J_X\) and the present principal-angle
threshold are equivalent up to subpower factors on that branch.
\end{proposition}

\begin{proof}
Use the exact dictionary
\[
 \rho_X=\canc_X\sqrt{\frac{\Seff(A_X)}{S_X}}
\]
from \cref{prop:tpc88-dictionary}.
\end{proof}

\subsection{Direct zero-mode error scale}

\begin{proposition}[Why \(J^{-1}\) produces a \(Q^2\) zero-mode
error]
\label{prop:direct-zero-mode-scale}
Suppose a same-frame surrogate has
\[
 S_X\ll X^{o(1)}Q_X,\qquad
 D(b_X)\ll X^{o(1)}Q_X^3J_X^2,
\]
and
\[
 \delta_X\le\frac{X^{o(1)}}{J_X}.
\]
Then
\begin{equation}
 |Z(a_X)-Z(b_X)|
 \le X^{o(1)}Q_X^2.
\label{eq:direct-zero-mode-scale}
\end{equation}
\end{proposition}

\begin{proof}
Apply \eqref{eq:zero-mode-transfer}:
\[
 |Z(a_X)-Z(b_X)|
 \le
 \sqrt{S_X}\,\delta_X\sqrt{D(b_X)}
 \le
 X^{o(1)}
 Q_X^{1/2}J_X^{-1}
 \bigl(Q_X^3J_X^2\bigr)^{1/2}
 =
 X^{o(1)}Q_X^2.
\]
\end{proof}

\begin{remark}
This computation explains the same threshold in the unnormalized
zero mode: \(J_X^{-1}\) is exactly the relative energy accuracy
that converts the natural determinant norm into the \(Q_X^2\)
zero-mode scale.  The proposition is conditional and provides no
such approximation for the physical coefficient.
\end{remark}
